\documentclass[border=12pt]{standalone}
\usepackage[no-math]{fontspec}
\usepackage[brazil]{babel}
\usepackage{amsmath,amssymb}
\IfFileExists{C:/Windows/Fonts/Lato_Regular.ttf}{%
  \setmainfont{Lato}[
    Path = C:/Windows/Fonts/,
    UprightFont = Lato_Regular.ttf,
    BoldFont = Lato_Bold.ttf,
    ItalicFont = Lato_Italic.ttf,
    BoldItalicFont = Lato_BoldItalic.ttf
  ]
  \setsansfont{Lato}[
    Path = C:/Windows/Fonts/,
    UprightFont = Lato_Regular.ttf,
    BoldFont = Lato_Bold.ttf,
    ItalicFont = Lato_Italic.ttf,
    BoldItalicFont = Lato_BoldItalic.ttf
  ]
}{%
  \RequirePackage[default]{lato}
}
\renewcommand{\familydefault}{\sfdefault}

\usepackage{xcolor}
\usepackage{tikz}
\usepackage{pgfplots}
\pgfplotsset{compat=1.18}
\usetikzlibrary{arrows.meta,positioning,shapes.geometric,fit,backgrounds,calc,decorations.pathreplacing,patterns}

% Paleta institucional PLANTERR / UEFS (alinhada ao estilo planterr-slides.css / 09a)
\definecolor{etNavy}{HTML}{1F4268}
\definecolor{etNavyDark}{HTML}{152E4A}
\definecolor{etGold}{HTML}{B38A3D}
\definecolor{etGoldDark}{HTML}{8C6927}
\definecolor{etTeal}{HTML}{0E7C7B}
\definecolor{etForest}{HTML}{2D6A4F}
\definecolor{etAlert}{HTML}{9E2A2B}
\definecolor{etInk}{HTML}{26313D}
\definecolor{etMuted}{HTML}{5B6573}
\definecolor{etBorder}{HTML}{D5DEE8}

\definecolor{bgNavy}{HTML}{EDF2F7}
\definecolor{bgGold}{HTML}{F8F3E9}
\definecolor{bgTeal}{HTML}{EBF6F5}
\definecolor{bgForest}{HTML}{EEF5F0}
\definecolor{bgAlert}{HTML}{FDF2F2}
\definecolor{bgCanvas}{HTML}{FFFFFF}

\tikzset{
  >=Stealth,
  etBox/.style={
    draw=#1, fill=white, rounded corners=5pt, line width=1.15pt,
    inner sep=7pt, align=center, font=\footnotesize\sffamily, text=etInk
  },
  etTint/.style={
    draw=#1, fill=#1!8!white, rounded corners=5pt, line width=1.15pt,
    inner sep=7pt, align=center, font=\footnotesize\sffamily, text=etInk
  },
  etBadge/.style={
    fill=#1, text=white, rounded corners=9pt, inner xsep=8pt, inner ysep=3pt,
    font=\scriptsize\bfseries\sffamily
  },
  etTitle/.style={font=\normalsize\bfseries\sffamily, text=etNavy, align=center},
  etSub/.style={font=\footnotesize\sffamily, text=etMuted, align=center},
  etArrow/.style={->, >=Stealth, line width=1.35pt},
  etDArrow/.style={<->, >=Stealth, line width=1.35pt}
}


\begin{document}
\begin{tikzpicture}[
  pill/.style={
    draw=#1!60!white, fill=white, rounded corners=4pt, line width=0.9pt,
    inner xsep=8pt, inner ysep=4.5pt, text width=4.1cm, align=center,
    font=\footnotesize\sffamily, text=etInk
  }
]

  % Título principal
  \node[etTitle] (title) at (0, 4.15)
    {Modelo Causal de Ecossistema Empreendedor (Stam, 2015; Stam \& van de Ven, 2021)};

  % Coluna 1: Condições Estruturais (Framework)
  \node[font=\footnotesize\bfseries\sffamily, text=etNavy, align=center] (fw_hdr) at (-6.35, 2.75)
    {Condições Estruturais\\{\scriptsize\normalfont\textcolor{etMuted}{(Framework)}}};
  \node[pill=etNavy, inner ysep=7.5pt] (fw1) at (-6.35, 1.65) {Instituições Formais};
  \node[pill=etNavy, inner ysep=7.5pt] (fw2) at (-6.35, 0.50) {Cultura Empreendedora};
  \node[pill=etNavy, inner ysep=7.5pt] (fw3) at (-6.35, -0.65) {Infraestrutura Física};
  \node[pill=etNavy, inner ysep=7.5pt] (fw4) at (-6.35, -1.80) {Demanda de Mercado};

  \begin{scope}[on background layer]
    \node[
      draw=etNavy, fill=bgNavy, rounded corners=6pt, line width=1.25pt,
      inner sep=9pt, fit=(fw_hdr)(fw1)(fw4)
    ] (framework) {};
  \end{scope}

  % Coluna 2: Condições Sistêmicas (Coração)
  \node[font=\footnotesize\bfseries\sffamily, text=etNavyDark, align=center] (sy_hdr) at (0, 2.75)
    {Condições Sistêmicas\\{\scriptsize\normalfont\textcolor{etMuted}{(Coração do Ecossistema)}}};
  \node[pill=etNavyDark] (sy1) at (0, 1.85) {Redes de Empreendedores};
  \node[pill=etNavyDark] (sy2) at (0, 1.12) {Liderança Territorial};
  \node[pill=etNavyDark] (sy3) at (0, 0.39) {Finanças e Capital};
  \node[pill=etNavyDark] (sy4) at (0, -0.34) {Talento e Qualificação};
  \node[pill=etNavyDark] (sy5) at (0, -1.07) {Conhecimento (P\&D)};
  \node[pill=etNavyDark] (sy6) at (0, -1.80) {Serviços Intermediários};

  \begin{scope}[on background layer]
    \node[
      draw=etNavyDark, fill=etNavy!11!white, rounded corners=6pt, line width=1.35pt,
      inner sep=9pt, fit=(sy_hdr)(sy1)(sy6)
    ] (systemic) {};
  \end{scope}

  % Coluna 3: Outputs e Outcomes empilhados
  \node[
    draw=etGold, fill=bgGold, rounded corners=6pt, line width=1.3pt,
    inner sep=10pt, align=center, text width=4.3cm, minimum height=2.0cm,
    font=\footnotesize\sffamily, text=etInk
  ] (outputs) at (6.45, 1.50) {
    \textbf{\textcolor{etGoldDark}{Outputs: Empreendedorismo\\Produtivo}}\\[5pt]
    {\scriptsize\textit{Startups e firmas escaláveis}}
  };

  \node[
    draw=etTeal, fill=bgTeal, rounded corners=6pt, line width=1.3pt,
    inner sep=10pt, align=center, text width=4.3cm, minimum height=2.0cm,
    font=\footnotesize\sffamily, text=etInk
  ] (outcomes) at (6.45, -1.25) {
    \textbf{\textcolor{etTeal}{Outcomes: Criação de Valor\\Agregado}}\\[5pt]
    {\scriptsize\textit{Desenvolvimento Regional}}
  };

  % Setas diretas (fluxo causal)
  \draw[etArrow, etNavy] (framework.east) -- (systemic.west);
  \draw[etArrow, etGold] (systemic.east |- outputs.west) -- (outputs.west);
  \draw[etArrow, etTeal] (outputs.south) -- (outcomes.north);

  % Seta de retroalimentação limpa ABAIXO de todas as caixas
  \coordinate (fb_y) at (0, -3.40);
  \draw[etArrow, etTeal, line width=1.45pt]
    (outcomes.south) -- (outcomes.south |- fb_y)
    -- node[
      fill=white, draw=etTeal, rounded corners=8pt, inner xsep=8pt, inner ysep=3pt,
      font=\scriptsize\bfseries\sffamily, text=etTeal
    ] {Retroalimentação (Reciclagem Empreendedora)}
    (systemic.south |- fb_y) -- (systemic.south);

\end{tikzpicture}
\end{document}
